\chapter{Introduction}
\label{cap:introducao}

Machine Learning (ML) models have become central to scientific research and
industrial applications, supporting decision-making processes in domains such as
healthcare, finance, engineering, and environmental sciences. As predictive
models are increasingly deployed in high-stakes settings, the reliable estimation
of their generalization performance becomes essential. Rigorous evaluation
protocols are required not only for comparing competing methods, but also to
ensure that reported results accurately reflect how models behave on previously
unseen data. However, when the process of training and validating ML models is
not conducted correctly from a methodological point of view, the reported
performance may diverge substantially from the model's real predictive
capability. \citeonline{kapoor2023} argue that such methodological issues are
widespread in ML-based scientific research, contributing to what has been
described as a reproducibility crisis in machine-learning-based science.

Among the factors that compromise evaluation reliability, data leakage (DL) is
one of the most critical and persistent. Data leakage is defined as the
contamination of the training process with privileged information from the
independent validation or test set, to which the algorithm should not have access
during model development \cite{kaufman2012}. Under a broader interpretation
\cite{kapoor2023}, leakage extends beyond direct train-test overlap and includes
more subtle methodological flaws, such as applying normalization using statistics
computed over the entire data set, imputing missing values before data splitting,
performing feature selection with information from both training and validation
sets, or tuning hyperparameters in a way that indirectly incorporates evaluation
data. Although these practices are frequently introduced unintentionally, they
bias performance estimates and may produce misleading conclusions about model
quality.

The relevance of the problem has motivated a growing body of research focused on
the characterization, detection, and mitigation of leakage, including formal
definitions \cite{kaufman2012}, taxonomies of leakage types \cite{kapoor2023},
static code analysis approaches \cite{drobnjakovic2022,yang2022}, and
domain-specific guidelines \cite{bernett2024,apicella2024}. Empirical
demonstrations of the phenomenon are comparatively scarcer, and the most
influential ones come from bioinformatics: \citeonline{ambroise2002} showed that
performing gene selection outside the cross-validation loop drives the estimated
error rate close to zero even on data with no real signal, and
\citeonline{varma2006} showed that the cross-validation error of a model whose
hyperparameters were themselves chosen by cross-validation is a significantly
biased estimate of the true error. These results, however, are confined to
classification and to a single application domain. More recently,
\citeonline{becker2023} conducted a systematic assessment of data leakage effects
on the performance estimates of ML classifiers across several data sets and
algorithms, showing that the magnitude of the distortion varies according to both
the learning algorithm and the pre-processing step in which the leakage
originates. Regression tasks, despite being equally susceptible to the same
methodological flaws, have not received an equivalent treatment, and it remains
unclear whether the patterns observed for classifiers transfer to continuous
prediction targets.

This work aims to address this gap by proposing a systematic and controlled study
of the effects of data leakage on the performance estimates of ML models
developed for regression tasks. The objective is to understand how different data
sets and learning algorithms are affected by each kind of pre-processing
contamination: normalization, missing values imputation, feature selection, and
hyperparameter optimization. For each of these steps, experiments contrast a
leakage-free pipeline with a leakage-prone counterpart under otherwise identical
conditions, using benchmark regression data sets obtained from the Penn Machine
Learning Benchmarks collection \cite{romano2022} and a diverse set of supervised
learning algorithms. By isolating each contamination mechanism, this work seeks
to quantify the magnitude and the variability of the performance inflation caused
by data leakage, and to investigate whether certain regression algorithms are
more sensitive to these distortions than others.

This work is organized as follows. Chapter~\ref{cap:fundamentacao} provides a
theoretical background, describing the key concepts for this work, including the
regression algorithms employed, the pre-processing and validation procedures
adopted, and the notion of data leakage. Chapter~\ref{cap:trabalhos} summarizes
existing efforts based on related works. Next, Chapter~\ref{cap:metodologia}
presents the methods used, including the selection of data and the experimental
methodology adopted for each contamination scenario. The experimental results are
unveiled in Chapter~\ref{cap:resultados}, alongside the research questions that
our work aimed to address. Finally, conclusions and possible future works are
presented in Chapter~\ref{cap:conclusao}.
